\documentclass{article}
\usepackage{amsmath,amssymb,amsthm}
\newtheorem{theorem}{Theorem}
\newtheorem{proposition}{Proposition}
\begin{document}
\title{Containment and signed-product selectors for finite planes}
\author{Joint-cell audit}
\date{24 November 2026}
\maketitle

\paragraph{Scope.}
The original xorformer model and its admissibility conditions are unchanged.
This note concerns a proposed implementation of the finite-plane incidence
selectors on larger Boolean row domains.  It does not give a lower bound
for arbitrary joint scores or for the original worst-case head complexity.
The first theorem and its energy identity are kernel checked in
\texttt{JointCellContainmentSelectors.lean}.  The subsequent observations
are written proofs only.

Let $F$ be a finite field of order $q$.  The points are $F^2$ and the
nonvertical lines are $L_{\alpha,\beta}=\{(x,y):y=\alpha x+\beta\}$.

\begin{theorem}[Containment dimension]
Suppose each point $p$ receives a set $T_p\subseteq[n]$ and each line $L$
receives a set $S_L\subseteq[n]$, with
\[
p\in L\quad\Longleftrightarrow\quad T_p\subseteq S_L.
\]
Then $n\geq q-1$.
\end{theorem}
\begin{proof}
For $f:F^2\to\mathbb R$, put $m_L=\sum_{p\in L}f(p)$.
The exact line-transform identity is
\[
\sum_L m_L^2
=q\sum_p f(p)^2+\left(\sum_p f(p)\right)^2
-\sum_x\left(\sum_y f(x,y)\right)^2.
\]
Indeed, a point belongs to $q$ lines, a pair with different first
coordinates belongs to one line, and two distinct points with the same
first coordinate belong to no nonvertical line.

For a point set $P$ and a line set $\mathcal L$ with no incidences, take
$f=1_P$ and $p=|P|$.  The first moment is $\sum_Lm_L=qp$ and the
identity gives $\sum_Lm_L^2\leq qp+p^2$.  Hence
\[
\sum_L(qm_L-p)^2\leq q^3p.
\]
Every line in $\mathcal L$ has $m_L=0$, so
$|\mathcal L|p^2\leq q^3p$.  Thus $|P||\mathcal L|\leq q^3$,
including the empty case.

For coordinate $i$, define
$P_i=\{p:i\in T_p\}$ and $\mathcal L_i=\{L:i\notin S_L\}$.
Their rectangle has no incidences.  Every nonincidence lies in at least
one such rectangle.  There are $q^4-q^3=q^3(q-1)$ nonincidences, and
therefore $q^3(q-1)\leq\sum_i|P_i||\mathcal L_i|\leq nq^3$.
\end{proof}

\paragraph{Pure-power consequence.}
On a $c$-hot Boolean address $1_T$, the squarefree coefficient of degree
$c$ in $(\sum_iw_i x_i)^c$ is $c!\prod_{i\in T}w_i$.
Its nonzero support is precisely $T\subseteq\{i:w_i\ne0\}$.
Consequently a single pure-power jet per line cannot encode a plane of
order $q>n+1$, regardless of how large $c$ is.  In particular increasing
the jet order alone does not create an exponentially large finite plane
inside only $n$ physical coordinates.

\paragraph{Two products with a positive ratio.}
There is a further restricted obstruction when all product factors are
nonzero and the ratio of the two products is positive on every address.
If a line selector is
\[
A_L\prod_{i\in T_p}u_{Li}-B_L\prod_{i\in T_p}v_{Li},
\]
then its zero set agrees with the zero set of
\[
\log|A_L/B_L|+\sum_{i\in T_p}\log|u_{Li}/v_{Li}|.
\]
The latter point-by-line matrix has real rank at most $n+1$.
If the former selector is nonzero exactly on line incidences, the latter
is a weighted incidence matrix with nonzero weight at every incidence.
The weighted design-matrix estimate from
\texttt{joint-cell-weighted-selectors.tex} then applies.
Opposite product signs can make the original expression nonzero while the
log ratio vanishes, so this real-rank argument alone is insufficient.
There is nevertheless a separate bound for arbitrary signs when all
factors remain nonzero.  Let $A$ be the real matrix of log absolute
ratios and let $B$ be the matrix over $\mathbb F_2$ recording whether
the two product signs differ.  Their ranks are at most $n+1$ over their
respective fields.  Both vanish off incidences and their supports jointly
cover every incidence.  Fix a parallel class of $q$ disjoint lines.  If
$r$ columns of $A$ are nonzero, those columns have disjoint supports and
are independent over $\mathbb R$.  The remaining $q-r$ columns of $B$
are nonzero and disjoint, hence independent over $\mathbb F_2$.
Therefore $r\leq n+1$ and $q-r\leq n+1$, giving $q\leq2(n+1)$.
This proof does not handle zero factors or three or more signed products.

\paragraph{Binary product addresses and a different rank test.}
Suppose a point $(x,y)\in F^2$ is encoded by separate binary encodings
of $x$ and $y$.  For every nonzero slope the $q$ by $q$ line-incidence
matrix is a permutation matrix.  A degree-$c$ polynomial of one affine
statistic $u(x)+v(y)$ has matrix rank at most $c+1$, by the binomial
expansion.  Thus a sum of $R$ such jets representing the line exactly
requires $R(c+1)\geq q$.
This statement depends on the product encoding and does not extend
without proof to arbitrary $c$-hot point codes.

\paragraph{What remains possible.}
Signed sums with several products, nonproduct point encodings, approximate
selectors combined with quantitative interference cancellation, and
joint rational decoders with independent column poles are not excluded
by these arguments.  Any successful construction still has to charge
its line selection, independent-memory corrections, proper-row support,
axis repair, and full-cube denominator positivity.
\end{document}
